\textbf{Phần 1. Trắc nghiệm lựa chọn}
\setcounterpageref{ex}{0}

\begin{ex} %[3D1V03]
Tìm giá trị của $x$, biết: $x + 145 = 320 + 215$.
\choice
{\True $390$}
{$380$}
{$400$}
{$410$}
\loigiai{
Tính giá trị của vế phải trước:\\
$$x + 145 = 535$$
$$x = 535 - 145 = 390$$
\textbf{Đáp án của câu hỏi là A}
}
\end{ex}

\begin{ex} %[3D1V03]
Tìm số bị trừ $x$ trong biểu thức: $x - (120 + 80) = 450$.
\choice
{$550$}
{$600$}
{\True $650$}
{$700$}
\loigiai{
Thực hiện phép tính trong ngoặc:\\
$$x - 200 = 450$$
$$x = 450 + 200 = 650$$
\textbf{Đáp án của câu hỏi là C}
}
\end{ex}

\begin{ex} %[3D1V03]
Tìm giá trị của $x$, biết: $850 - x = 240 + 160$.
\choice
{$400$}
{\True $450$}
{$500$}
{$350$}
\loigiai{
Tính vế phải:\\
$$850 - x = 400$$
$$x = 850 - 400 = 450$$
\textbf{Đáp án của câu hỏi là B}
}
\end{ex}

\begin{ex} %[3D1V03]
Tổng của hai số bằng $480$. Nếu giữ nguyên số hạng thứ nhất và tăng số hạng thứ hai thêm $45$ đơn vị thì tổng mới là bao nhiêu?
\choice
{$515$}
{$535$}
{$435$}
{\True $525$}
\loigiai{
Khi giữ nguyên một số hạng và tăng số hạng kia thêm bao nhiêu đơn vị thì tổng cũng tăng thêm bấy nhiêu đơn vị.\\
Tổng mới là:
$$480 + 45 = 525$$
\textbf{Đáp án của câu hỏi là D}
}
\end{ex}

\begin{ex} %[3D1V03]
Hai số có tổng bằng $750$. Biết rằng nếu bớt ở số thứ nhất đi $50$ đơn vị thì hai số bằng nhau. Số thứ nhất ban đầu là
\choice
{\True $400$}
{$350$}
{$450$}
{$500$}
\loigiai{
Vì bớt ở số thứ nhất đi $50$ đơn vị thì hai số bằng nhau nên số thứ nhất hơn số thứ hai $50$ đơn vị.\\
Hai lần số thứ hai là:
$$750 - 50 = 700$$
Số thứ hai là:
$$700 : 2 = 350$$
Số thứ nhất ban đầu là:
$$350 + 50 = 400$$
\textbf{Đáp án của câu hỏi là A}
}
\end{ex}

\begin{ex} %[3D1V03]
Một phép trừ có số bị trừ bằng $540$. Nếu tăng số trừ thêm $30$ đơn vị thì được hiệu mới bằng $180$. Số trừ ban đầu là
\choice
{$300$}
{$360$}
{\True $330$}
{$320$}
\loigiai{
Khi số trừ tăng thêm $30$ đơn vị thì hiệu giảm đi $30$ đơn vị.\\
Hiệu ban đầu của phép trừ là:
$$180 + 30 = 210$$
Số trừ ban đầu là:
$$540 - 210 = 330$$
\textbf{Đáp án của câu hỏi là C}
}
\end{ex}

\begin{ex} %[3D1V03]
Tìm giá trị của $x$, biết: $900 - x - 150 = 420$.
\choice
{$320$}
{\True $330$}
{$340$}
{$350$}
\loigiai{
Biến đổi biểu thức:\\
$$900 - 150 - x = 420$$
$$750 - x = 420$$
$$x = 750 - 420 = 330$$
\textbf{Đáp án của câu hỏi là B}
}
\end{ex}

\begin{ex} %[3D1V03]
Trong một phép trừ, số trừ là $135$ và hiệu lớn hơn số trừ $40$ đơn vị. Số bị trừ là
\choice
{$300$}
{$290$}
{$320$}
{\True $310$}
\loigiai{
Hiệu của phép trừ là:
$$135 + 40 = 175$$
Số bị trừ là:
$$175 + 135 = 310$$
\textbf{Đáp án của câu hỏi là D}
}
\end{ex}

\begin{ex} %[3D1V03]
Bác An có hai bao thóc nặng tất cả $180\text{ kg}$. Nếu bác An chuyển từ bao thứ nhất sang bao thứ hai $15\text{ kg}$ thì hai bao nặng bằng nhau. Bao thứ nhất lúc đầu nặng số ki-lô-gam là
\choice
{\True $105\text{ kg}$}
{$95\text{ kg}$}
{$100\text{ kg}$}
{$110\text{ kg}$}
\loigiai{
Khi hai bao nặng bằng nhau thì mỗi bao nặng số ki-lô-gam là:
$$180 : 2 = 90\text{ (kg)}$$
Bao thứ nhất lúc đầu nặng số ki-lô-gam là:
$$90 + 15 = 105\text{ (kg)}$$
\textbf{Đáp án của câu hỏi là A}
}
\end{ex}

\begin{ex} %[3D1V03]
Tìm chữ số $x$ biết: $\overline{x45} + 230 = 875$.
\choice
{$5$}
{$7$}
{\True $6$}
{$8$}
\loigiai{
Ta có:
$$\overline{x45} = 875 - 230 = 645$$
Do đó chữ số $x = 6$.\\
\textbf{Đáp án của câu hỏi là C}
}
\end{ex}

\begin{ex} %[3D1V03]
Một kho có một số ki-lô-gam gạo. Sau khi xuất đi $160\text{ kg}$ gạo rồi lại nhập thêm $90\text{ kg}$ gạo thì trong kho có tất cả $350\text{ kg}$ gạo. Lúc đầu trong kho có số ki-lô-gam gạo là
\choice
{$400\text{ kg}$}
{\True $420\text{ kg}$}
{$440\text{ kg}$}
{$380\text{ kg}$}
\loigiai{
Gọi số gạo ban đầu trong kho là $x\text{ (kg)}$. Theo đề bài ta có:\\
$$x - 160 + 90 = 350$$
$$x - 70 = 350$$
$$x = 350 + 70 = 420\text{ (kg)}$$
\textbf{Đáp án của câu hỏi là B}
}
\end{ex}

\begin{ex} %[3D1V03]
Tìm giá trị của $x$, biết: $(x + 210) - 150 = 380 + 120$.
\choice
{$420$}
{$430$}
{$450$}
{\True $440$}
\loigiai{
Tính vế phải:\\
$$(x + 210) - 150 = 500$$
$$x + 210 = 500 + 150$$
$$x + 210 = 650$$
$$x = 650 - 210 = 440$$
\textbf{Đáp án của câu hỏi là D}
}
\end{ex}

\textbf{Phần 2. Trắc nghiệm đúng sai}
\setcounterpageref{ex}{0}

\begin{ex} %[3D1V03]
Xét tính đúng sai của các nhận định sau về các quy tắc tìm thành phần phép tính:
\choiceTF
{\True Muốn tìm một số hạng chưa biết trong một tổng, ta lấy tổng trừ đi số hạng đã biết.}
{Muốn tìm số trừ, ta lấy số bị trừ cộng với hiệu.}
{\True Nếu $x - 50 = 200$ thì giá trị của $x$ là $250$.}
{\True Trong một phép trừ, nếu hiệu bằng $0$ thì số bị trừ bằng số trừ.}
\loigiai{
\begin{itemize}
    \item Ý a) Đúng: Theo quy tắc tìm số hạng chưa biết.
    \item Ý b) Sai: Muốn tìm số trừ, ta phải lấy số bị trừ trừ đi hiệu ($\text{Số trừ} = \text{Số bị trừ} - \text{Hiệu}$).
    \item Ý c) Đúng: $x = 200 + 50 = 250$.
    \item Ý d) Đúng: Khi $\text{Số bị trừ} - \text{Số trừ} = 0 \Rightarrow \text{Số bị trừ} = \text{Số trừ}$.
\end{itemize}
\begin{itemize}
    \item a) \textbf{Đúng}.
    \item b) \textbf{Sai}.
    \item c) \textbf{Đúng}.
    \item d) \textbf{Đúng}.
\end{itemize}
}
\end{ex}

\begin{ex} %[3D1V03]
Cho bài toán tìm $x$: $(500 - x) + 120 = 320$. Xét tính đúng sai của các khẳng định sau:
\choiceTF
{\True Giá trị của biểu thức $(500 - x)$ là $200$.}
{Giá trị của $x$ tìm được là $700$.}
{\True Nếu vế phải tăng thêm $100$ đơn vị (thành $420$) thì giá trị của $x$ sẽ giảm đi $100$ đơn vị.}
{\True Số $x$ tìm được là một số tròn trăm lớn hơn $250$.}
\loigiai{
\begin{itemize}
    \item Ý a) Đúng: $500 - x = 320 - 120 = 200$.
    \item Ý b) Sai: $x = 500 - 200 = 300$ (không phải 700).
    \item Ý c) Đúng: Khi vế phải là 420 thì $500 - x = 420 - 120 = 300 \Rightarrow x = 500 - 300 = 200$. Giá trị $x$ giảm từ $300$ xuống $200$ (giảm 100 đơn vị).
    \item Ý d) Đúng: $x = 300$ là số tròn trăm và $300 > 250$.
\end{itemize}
\begin{itemize}
    \item a) \textbf{Đúng}.
    \item b) \textbf{Sai}.
    \item c) \textbf{Đúng}.
    \item d) \textbf{Đúng}.
\end{itemize}
}
\end{ex}

\textbf{Phần 3. Trả lời ngắn}
\setcounterpageref{ex}{0}

\begin{ex} %[3D1V03]
Tìm giá trị của $x$, biết: $x + 256 = 600 - 144$.
\shortans{200}
\loigiai{
\textbf{Đáp số:} $200$\\
Tính vế phải:\\
$$x + 256 = 456$$
$$x = 456 - 256 = 200$$
}
\end{ex}

\begin{ex} %[3D1V03]
Tìm một số tự nhiên, biết rằng nếu lấy số đó trừ đi $170$ rồi cộng thêm $90$ thì được kết quả là $420$.
\shortans{500}
\loigiai{
\textbf{Đáp số:} $500$\\
Gọi số cần tìm là $x$. Ta có:\\
$$x - 170 + 90 = 420$$
$$x - 80 = 420$$
$$x = 420 + 80 = 500$$
}
\end{ex}

\begin{ex} %[3D1V03]
Hai số có hiệu bằng $120$. Nếu thêm vào số bị trừ $25$ đơn vị và giữ nguyên số trừ thì hiệu mới bằng bao nhiêu?
\shortans{145}
\loigiai{
\textbf{Đáp số:} $145$\\
Khi số bị trừ tăng thêm bao nhiêu đơn vị và giữ nguyên số trừ thì hiệu cũng tăng thêm bấy nhiêu đơn vị.\\
Hiệu mới là:
$$120 + 25 = 145$$
}
\end{ex}

\begin{ex} %[3D1V03]
Một bác nông dân mang một số tiền đi chợ. Bác mua một bao phân bón hết $65$ nghìn đồng, mua giống cây trồng hết $45$ nghìn đồng thì bác còn lại đúng $90$ nghìn đồng. Hỏi ban đầu bác nông dân có tất cả bao nhiêu nghìn đồng?
\shortans{200}
\loigiai{
\textbf{Đáp số:} $200$\\
Số tiền bác nông dân đã tiêu là:
$$65 + 45 = 110\text{ (nghìn đồng)}$$
Ban đầu bác nông dân có số tiền là:
$$110 + 90 = 200\text{ (nghìn đồng)}$$
}
\end{ex}

\begin{ex} %[3D1V03]
Tìm giá trị của $x$, biết: $780 - (x + 130) = 250$.
\shortans{400}
\loigiai{
\textbf{Đáp số:} $400$\\
Xem $(x + 130)$ là số trừ chưa biết:\\
$$x + 130 = 780 - 250$$
$$x + 130 = 530$$
$$x = 530 - 130 = 400$$
}
\end{ex}

\begin{ex} %[3D1V03]
Trong một phép cộng có hai số hạng, nếu tăng số hạng thứ nhất thêm $48$ đơn vị và giảm số hạng thứ hai đi $23$ đơn vị thì được tổng mới bằng $365$. Tìm tổng ban đầu của hai số.
\shortans{340}
\loigiai{
\textbf{Đáp số:} $340$\\
Vì tăng số hạng thứ nhất thêm 48 đơn vị và giảm số hạng thứ hai đi 23 đơn vị nên tổng mới tăng thêm số đơn vị là:
$$48 - 23 = 25\text{ (đơn vị)}$$
Tổng ban đầu của hai số là:
$$365 - 25 = 340$$
}
\end{ex}
